\section*{September 8, 2026: What Huron's two floor areas measure}
\addcontentsline{toc}{section}{September 8, 2026: What Huron's two floor areas measure}

At 910 W Huron, the increase from 176,608 to 295,452 square feet is an exact change in the scope of the Assessor record. In 2021, one tax parcel contains all 141 apartments and 176,608 square feet. A second parcel at the same address contains zero apartments and 118,844 square feet, also recorded as commercial space. The 2024 record combines the two parcels and their floor areas. This explains the arithmetic without establishing whether the second area belongs in the paper's floor-area measure.

The second parcel's public building cards describe air rights but leave their floor areas blank. A parking structure is present at the building, but the inspected records do not establish that all 118,844 square feet are parking or that the two measurements do not overlap. The City's energy records do not resolve this: they report 171,344 square feet in 2016--2019, 300,047 in 2020--2021, and 177,566 in 2022. The larger planned development's zoning limits cannot be assigned to this building alone.

Jacob requested further investigation before any exclusion. No new Huron floor-area value or eligibility decision has been adopted. The remaining question requires a building-area schedule or assessment worksheet that identifies the spaces in the second parcel. The frozen paper input and regressions are unchanged.

\smallskip
\noindent\textit{Reproduction note (Codex draft):} The commercial comparison uses the pinned 2021 and 2024 valuation records for parcel numbers ending 213-009 and 213-011. The September 8 City energy extract and PD 356 documents are acquired through the development-record acquisition task; checksums and a standard energy-data report record the sources. The release-audit land/location memo records the web evidence and remaining uncertainty. This entry documents an investigation, not a completed data correction.
